% =============================================================
%  第二单元  测量（一）
% =============================================================
\Chapter{第二关\quad 尺子、绳子和一双大手}
\noindent{\small\kaishu\color{tipc}（课本第二单元：测量（一），P20--P29）}

\begin{blunder}
去年网上买床单，豆角爸爸把“床宽”记成了“床长”，买回来短了 20 厘米。
铺上去像条围裙，四边都露着床垫。小豆角说：“豆角爸爸，这床单是不是缩水了？”
豆角爸爸没敢说是自己量错的。
\end{blunder}

\begin{gamble}{赌一把：估长度，谁准谁赢}
\begin{itemize}
  \item 随便指家里一样东西（沙发、电视、门），两个人\textbf{各写一个估计数}，不许说出口。
  \item 一起量。\textbf{谁差得少谁赢}。输的人请对方一瓶小可乐，
        \textbf{但自己一口都不能喝}。连输三天换人。
  \item 进阶：不许用尺子，只能用“爸爸的一拃”去量。这时候小豆角会发现——
        大人的一拃和小朋友的一拃，根本不是一回事。
\end{itemize}
\end{gamble}

\begin{bonus}
\begin{itemize}
  \item \textbf{身高比赛}：在门框上画线，两个人都量，做一张榜贴墙上。
        每月更新一次，让小豆角自己读数、自己记日期。\textbf{小豆角会天天去看自己长高没有。}
  \item \textbf{买鞋量脚}：让小豆角自己把脚画在纸上，量最长的那根指头到脚跟，
        记下多少厘米，对照鞋码表看对不对。
  \item \textbf{不用尺子量床}：只许用小豆角的“一拃”去量小豆角的床，报“几拃”。
        然后豆角爸爸用手量一遍，再比一比——为什么数不一样？（这一关就通了）
\end{itemize}
\end{bonus}

\begin{unitinfo}
\textbf{课本 P20--P29}\quad \textbf{5 次课，每次 20 分钟}\quad
\textbf{建议第 5--7 周} \\[4pt]
\textbf{这一关到底在干嘛}：先让小豆角吃一次“量出来不算数”的亏，
再教小豆角把尺子用对。\\
\textbf{逻辑是这么走的}：\ 大家各量各的 $\to$ 结果对不上 $\to$ 所以要有个统一的“一小段”
$\to$ 这一小段叫 1 厘米 $\to$ 10 个 10 厘米拼起来叫 1 米。\\
\textbf{小豆角最容易栽的地方}：不是读数，是\textbf{没有单位感}。
能读对尺子，却把铅笔填成“17 米”。
\end{unitinfo}

\subsection{开局先聊个事：这些事生活里天天遇到}

给小豆角买鞋，得量脚长多少厘米；网上买床单，得看床多宽；
带小豆角做核酸，脚下一米线——都是一样的东西。

至于“17 米长的铅笔”这种笑话，其实大人也常犯：
有人问豆角爸爸客厅多大，豆角爸爸说“二十来平”，问豆角爸爸多远，豆角爸爸说“走过去五分钟”。
\textbf{单位感}这件事，靠的是身体记住，不是靠背。

所以这一章分两步：\textbf{先让小豆角量乱一次，再教小豆角量准。}

\subsection{任务 1：故意让小豆角量乱一次}

\begin{lesson}{1}{20}
\textbf{今天要拿下的}：小豆角自己得出“必须用一样长的东西量，别人才听得懂”。

\begin{materials}
一张桌子（或书桌）、铅笔、筷子。\textbf{今天不给小豆角尺子。}
\end{materials}

\begin{enumerate}
  \item \textbf{给任务}：量一量这张桌子有多长。\textbf{不许用尺子。}
  \item \textbf{第一轮用手}：用“拃”（张开大拇指到中指）量，边量边数。
        记下来：“桌子有 \blank 拃长”。
  \item \textbf{第二轮用铅笔}：再量。记：“\blank 支铅笔长”。
  \item \textbf{第三轮用筷子}：再记一次。
  \item \textbf{三个数摆一起}，问：“同一张桌子，怎么三个数？”
        ——小豆角多半会说“因为东西长短不一样”。
  \item \textbf{换人来（关键）}：让小豆角量完，豆角爸爸再用手量一遍。数又不一样了（豆角爸爸手大）。
        问：“豆角爸爸说 6 拃，小豆角说 8 拃，谁对？”
        小豆角会发现：\textbf{都对，但没法交流。}
  \item \textbf{让小豆角自己下结论}：“要用\textbf{一样长的东西}量，别人才听得懂。”
        这时候豆角爸爸才接一句：“所以全世界约好了一小段长度，叫厘米。”
  \item \textbf{顺便让小豆角总结“量的时候要注意什么”}：摆直、不能有空隙、每量一次做记号。
        这三条小豆角自己说出来，比讲十遍管用。
\end{enumerate}

\begin{saybox}
“同一张桌子，怎么三个数字？”\\
“小豆角说 8 拃，豆角爸爸说 6 拃，谁对？”（都对 $\to$ 所以要统一）\\
“小豆角的手歪了一点，量出来会变大还是变小？”\\
“上一次量到哪儿，下一次从哪儿开始？小豆角怎么记住？”（引出做记号）
\end{saybox}

\begin{passbox}
\begin{itemize}
  \item 小豆角能说出“要用一样长的东西量，才好比”。
  \item 小豆角自己提出了“要做记号”“要摆直”这两条。
\end{itemize}
\end{passbox}
\end{lesson}

\subsection{任务 2：尺子怎么拿才对}

\begin{lesson}{2}{20}
\textbf{今天要拿下的}：0 刻度对齐、贴平、眼睛正着看；手上有一厘米的感觉。

\begin{materials}
15\,cm 透明直尺两把、8 样要量的东西（橡皮、铅笔、文具盒、课本、手指、筷子……）。
\end{materials}

\begin{enumerate}
  \item \textbf{认尺子}：让小豆角自己找 \textbf{0 刻度}——是尺子上那条写着 0 的线，
        \textbf{不是尺子最边上那条边}。再找 1 厘米有多长，再找“cm”两个字母。
  \item \textbf{一厘米的手感}：食指和拇指比出 1 厘米，保持住，再看尺子对不对。
        \textbf{再找身上三个 1 厘米}：食指指甲盖宽度、门牙宽度、耳朵厚度。
  \item \textbf{三条规矩}（边说边做）：
        \begin{itemize}
          \item \textbf{0 对齐}：东西一头对准 0（不是对准尺子边）。
          \item \textbf{贴平}：尺子贴着东西，跟被量的那条边平行。
          \item \textbf{正着看}：眼睛对着读数，歪着看会读错。
        \end{itemize}
  \item \textbf{让小豆角判对错}：下面三种量法，哪个对？
        \begin{center}
        \begin{tikzpicture}[x=0.36cm,y=0.5cm,font=\scriptsize]
          \draw[thick] (0,0) rectangle (11,0.9);
          \foreach \i in {0,...,10}{ \draw (\i,0)--(\i,0.45); \node[below,inner sep=1pt] at (\i,0) {\i}; }
          \node[right,inner sep=2pt] at (11,0.45) {cm};
          \draw[matc,line width=3pt] (0,1.05)--(6,1.05);
          \node[above,matc] at (3,1.05) {一支笔：0 对齐（对）};
          \begin{scope}[yshift=-2.1cm]
          \draw[thick] (-0.5,0) rectangle (10.5,0.9);
          \foreach \i in {0,...,10}{ \draw (\i,0)--(\i,0.45); \node[below,inner sep=1pt] at (\i,0) {\i}; }
          \draw[pitc,line width=3pt] (-0.5,1.05)--(5.5,1.05);
          \node[above,pitc] at (2.5,1.05) {从尺子边缘量（错，会多 1 厘米）};
          \end{scope}
          \begin{scope}[yshift=-4.2cm]
          \draw[thick] (0,0) rectangle (11,0.9);
          \foreach \i in {0,...,10}{ \draw (\i,0)--(\i,0.45); \node[below,inner sep=1pt] at (\i,0) {\i}; }
          \draw[pitc,line width=3pt] (2,1.05)--(7,1.05);
          \node[above,pitc] at (4.5,1.05) {从 2 量到 7（也行，但要记住起点）};
          \end{scope}
        \end{tikzpicture}
        \end{center}
        给小豆角的口诀是：“\textbf{从 0 开始最省事；不从 0，就用大数减小数。}”
  \item \textbf{量 8 样东西}（填下表），\textbf{每样先估再量，两栏都填}：
        \begin{center}
        \small
        \begin{tabular}{@{}p{3.2cm}p{2.6cm}p{2.6cm}p{3.4cm}@{}}
        \toprule
        \textbf{量什么} & \textbf{豆角爸爸先估} & \textbf{量出来} & \textbf{差多少} \\
        \midrule
        橡皮长 & \blank[1.6cm] & \blank[1.6cm] & \blank[1.6cm] \\
        铅笔长 & \blank[1.6cm] & \blank[1.6cm] & \blank[1.6cm] \\
        文具盒长 & \blank[1.6cm] & \blank[1.6cm] & \blank[1.6cm] \\
        课本的长 & \blank[1.6cm] & \blank[1.6cm] & \blank[1.6cm] \\
        课本的宽 & \blank[1.6cm] & \blank[1.6cm] & \blank[1.6cm] \\
        小豆角的手指长 & \blank[1.6cm] & \blank[1.6cm] & \blank[1.6cm] \\
        小豆角的手掌宽 & \blank[1.6cm] & \blank[1.6cm] & \blank[1.6cm] \\
        筷子长 & \blank[1.6cm] & \blank[1.6cm] & \blank[1.6cm] \\
        \bottomrule
        \end{tabular}
        \end{center}
  \item \textbf{看最后一栏}：哪几样估得准？为什么？（天天摸的东西，估得准。
        所以“估”不是天赋，是经验。）
\end{enumerate}

\begin{saybox}
“0 在哪儿？是尺子最边上那条线吗？”\\
“笔的一头对在 0，另一头对着几？”\\
“铅笔大概 17——这个 17 是什么单位？”（让小豆角自己说出“厘米”两个字）\\
“小豆角估的和量的差了几厘米？下次怎么估得更准？”
\end{saybox}

\begin{pitbox}
\textbf{三个高频错误}：
\begin{enumerate}[label=\arabic*、]
  \item \textbf{从尺子边缘量}（把“尺子边”当成 0）。\textbf{解法}：让小豆角用手指点住 0，再把笔放上去。
  \item \textbf{不从 0 又不换算}：从 3 量到 8，直接读 8 厘米。
        \textbf{解法}：“小豆角数了几个格？5 个格就是 5 厘米，$8-3=5$。”
  \item \textbf{只写数字不写单位}。\textbf{解法}：整章必须写“5 厘米”，不许只写 5。
\end{enumerate}
\end{pitbox}

\begin{passbox}
\begin{itemize}
  \item 8 样东西独立量完，误差不超过 1 厘米，单位全写。
  \item 小豆角能说出量长度的三个要点。
\end{itemize}
\end{passbox}
\end{lesson}

\subsection{任务 3：画线段，还有“坏尺子”}

\begin{lesson}{3}{20}
\textbf{今天要拿下的}：会画指定长度的线段；尺子的 0 看不见了也能画。

\begin{materials}
15\,cm 直尺、一把 \textbf{“坏尺子”}（用胶带把 0 那头贴住）、纸、笔。
\end{materials}

\begin{enumerate}
  \item \textbf{画 4 厘米的线，就三步}：
        ① 笔尖在 0 刻度和纸上点一个点；② 在 4 刻度也点一个点；
        ③ 尺子连起来，\textbf{并写上“4 厘米”}。
        \textbf{写完长度才算画完}——这是规矩，考试也这么要求。
  \item \textbf{再画 8 厘米}，让小豆角自己说那三步。
  \item \textbf{照样子画一条}：豆角爸爸先画一条 6 厘米的，小豆角画一条一样长的。
  \item \textbf{坏尺子挑战（本课重点）}：拿出贴住 0 的尺子。
        “这把尺子废了，0 看不见，小豆角能画 6 厘米吗？”
        \begin{center}
        \begin{tikzpicture}[x=0.55cm,y=0.62cm,font=\scriptsize]
          \draw[thick] (0,0) rectangle (8,1);
          \fill[gray!30] (0,0) rectangle (3,1);
          \node at (1.5,0.5) {0 看不见};
          \foreach \i in {3,...,8}{ \draw (\i,0)--(\i,0.45);
            \node[below,inner sep=2pt] at (\i,0) {\i}; }
          \node[above right,inner sep=1pt] at (8,1) {cm};
        \end{tikzpicture}
        \end{center}
        引导：\textbf{从 3 开始，$3+6=9$，画到 9}。
        也可以从 4 到 10、从 5 到 11——\textbf{只要两个刻度差 6 就行}。
  \item \textbf{让小豆角试两个不同起点}，体会“起点不同，长度一样”。
  \item \textbf{反过来}：给一条已经画好的线（从 3 到 8），问多长？$8-3=5$ 厘米。
  \item \textbf{收一句技巧}：“不从 0 开始画，就记住\textbf{起点 + 长度 = 终点}。”
\end{enumerate}

\begin{saybox}
“0 看不见了，小豆角想从哪条线开始？从那儿画多长能到 9？”\\
“从 3 画到 9 是 6 厘米，从 4 画到 10 也是 6 厘米，为什么？”\\
“小豆角这条线不写‘6 厘米’，别人知道多长吗？”
\end{saybox}

\begin{passbox}
\begin{itemize}
  \item 4 cm、8 cm 画得规范，长度都标上。
  \item \textbf{能用坏尺子画出 6 厘米，并说清道理。}
  \item 能算出“从 3 画到 8”是 5 厘米。
\end{itemize}
\end{passbox}
\end{lesson}

\subsection{任务 4：1 米到底多长}

\begin{lesson}{4}{20}
\textbf{今天要拿下的}：身上有一米的感觉；记住 $1\ \text{米}=100\ \text{厘米}$；
会写成“几米几厘米”。

\begin{materials}
1 米长的绳子（提前剪好）、软尺、卷尺。
\end{materials}

\begin{enumerate}
  \item \textbf{拉直绳子}：两手一拉，这就是 1 米。
        一般比小豆角的两臂侧平举短一点点。
  \item \textbf{往身上比}：
        \begin{itemize}
          \item 让小豆角靠墙站好，豆角爸爸量小豆角身高（二年级一般 120--135 厘米）。
          \item “1 米到小豆角身体哪儿？”（大概胸口往上、肩膀附近）
          \item “小豆角大概几个 1 米高？”（一个多）
        \end{itemize}
  \item \textbf{去比家里的东西}，比 1 米长的打“√”，短的打“×”：
        \begin{center}
        \small
        \begin{tabular}{@{}p{3.4cm}p{1.2cm}p{3.4cm}p{1.2cm}@{}}
         门高 \blank[1cm] & & 餐桌长 \blank[1cm] & \\
         沙发宽 \blank[1cm] & & 书桌高 \blank[1cm] & \\
         小豆角的床长 \blank[1cm] & & 衣柜宽 \blank[1cm] & \\
        \end{tabular}
        \end{center}
  \item \textbf{把换算做实}：让小豆角在纸上写一遍 \ $1\ \text{米} = 100\ \text{厘米}$，
        然后数尺子：10、20、30……100，一共 10 个 10 厘米。
        \textbf{“10 个 10 厘米才是 1 米”，不是 10 厘米。}这一条必须亲手数一次。
  \item \textbf{量黑板（或床）}：用 1 米绳子一段一段量，剩下的零头用尺子量。
        记成“\blank 米 \blank 厘米”。
        \textbf{这是小豆角最容易丢零头的地方}——一定要小豆角把零头写出来。
  \item \textbf{口答换算}：$1\ \text{米}=$\blank\,厘米；$200\ \text{厘米}=$\blank\,米；
        $1\ \text{米}20\ \text{厘米}=$\blank\,厘米；
        $150\ \text{厘米}=$\blank\,米\blank\,厘米。
\end{enumerate}

\begin{saybox}
“1 米到小豆角哪儿？比给豆角爸爸看。”\\
“黑板量了两整段，还多出 35 厘米，那一共多长？”（2 米 35 厘米）\\
“300 厘米是几米？小豆角怎么想的？”（100 厘米是 1 米，3 个就是 3 米）
\end{saybox}

\begin{pitbox}
\textbf{两处最容易错}：
\begin{enumerate}[label=\arabic*、]
  \item \textbf{零头丢掉}：量出“2 米多一点”，只写 2 米。
        \textbf{解法}：必须写成“2 米 35 厘米”，零头不许省。
  \item \textbf{把 1 米记成 10 厘米}。\textbf{解法}：让小豆角自己在尺子上数到 100，
        看到底几个 10。
\end{enumerate}
\end{pitbox}

\begin{passbox}
\begin{itemize}
  \item 能用手比出 1 米，说清 1 米到自己身上哪儿。
  \item $1\ \text{米}=100\ \text{厘米}$ 正反都对，包括“米+厘米”的混合。
  \item 能量出床的长度，写成“X 米 Y 厘米”。
\end{itemize}
\end{passbox}
\end{lesson}

\subsection{任务 5：估得准，才算真会}

\begin{lesson}{5}{20}
\textbf{今天要拿下的}：会给东西挑对单位（厘米还是米）；越估越准。

\begin{materials}
软尺、笔、这个表格。
\end{materials}

\begin{enumerate}
  \item \textbf{单位体检}（快速过一遍，说“厘米”还是“米”）：
        \begin{center}
        \small
        \begin{tabular}{@{}ll@{}}
        门高约 2 \blank[1.6cm] & 黄瓜长约 20 \blank[1.6cm] \\
        李老师身高约 175 \blank[1.6cm] & 一拃长约 15 \blank[1.6cm] \\
        橡皮长约 3 \blank[1.6cm] & 树高约 10 \blank[1.6cm] \\
        铅笔长约 17 \blank[1.6cm] & 旗杆高约 8 \blank[1.6cm] \\
        房子高约 3 \blank[1.6cm] & 毛巾长约 60 \blank[1.6cm] \\
        杯子高约 9 \blank[1.6cm] & 粉笔长约 7 \blank[1.6cm] \\
        \end{tabular}
        \end{center}
        答案：米、厘米、厘米、厘米、厘米、米、厘米、米、米、厘米、厘米、厘米。
  \item \textbf{给小豆角判断办法}（不要死记）：
        “\textbf{拿在手里的、比手臂短的小东西，用厘米；比 1 米还长的大东西，用米。}”
        再补一句身体尺：桌高约 75 厘米（不到 1 米），教室长约 8 米（比手长太多）。
  \item \textbf{估一估、量一量（做两轮）}：
        \begin{center}
        \small
        \begin{tabular}{@{}p{3.2cm}p{2.6cm}p{2.6cm}p{3.4cm}@{}}
        \toprule
        \textbf{东西} & \textbf{第一次估} & \textbf{量出来} & \textbf{差多少} \\
        \midrule
        铅笔长 & \blank[1.6cm] & \blank[1.6cm] & \blank[1.6cm] \\
        文具盒长 & \blank[1.6cm] & \blank[1.6cm] & \blank[1.6cm] \\
        餐桌长 & \blank[1.6cm] & \blank[1.6cm] & \blank[1.6cm] \\
        门高 & \blank[1.6cm] & \blank[1.6cm] & \blank[1.6cm] \\
        小豆角的身高 & \blank[1.6cm] & \blank[1.6cm] & \blank[1.6cm] \\
        \bottomrule
        \end{tabular}
        \end{center}
        \textbf{重点看最后一栏}：第一次差 20 厘米，第二次只差 3 厘米——
        这就是“估的本事长出来了”。
  \item \textbf{立定跳远（动起来，课本 P27）}：连跳 4 次，每次量，记下来。
        \begin{center}
        \small
        \begin{tabular}{@{}p{2.6cm}p{1.8cm}p{1.8cm}p{1.8cm}p{1.8cm}@{}}
        \toprule
        & 第一次 & 第二次 & 第三次 & 第四次 \\
        \midrule
        成绩（厘米） & \blank[1.2cm] & \blank[1.2cm] & \blank[1.2cm] & \blank[1.2cm] \\
        \bottomrule
        \end{tabular}
        \end{center}
        问：最远多少？最近多少？最远比最近多多少？（顺手复习第一单元的减法）
  \item \textbf{加点料}（课本 P29）：中国自己的长度单位——
        \textbf{1 米 = 3 尺，1 尺 = 10 寸}。
        让小豆角用手比一比 1 尺（约 33 厘米）。
        可以问：“古人说‘七尺男儿’，大概多高？”（古尺更短，别深究，笑一下就行）
\end{enumerate}

\begin{saybox}
“这根黄瓜 20 米还是 20 厘米？20 米大概多长？”（比到楼下——小豆角一听就笑，
笑声就是单位感）\\
“小豆角怎么知道门是 2 米不是 2 厘米？”\\
“小豆角第二次估准多了，小豆角用了什么办法？”（把小豆角的答案变成小豆角的方法）
\end{saybox}

\begin{passbox}
\begin{itemize}
  \item 12 道选单位全对。
  \item 小豆角能说出“怎么判断”，不是靠猜。
  \item 估的误差从“离谱”缩到 10 厘米以内。
\end{itemize}
\end{passbox}
\end{lesson}

\subsection{翻车现场（对着看）}

\begin{center}
\small
\begin{tabularx}{\linewidth}{@{}p{3.8cm}p{4.0cm}X@{}}
\toprule
\textbf{小豆角写出来的样子} & \textbf{其实是想错了什么} & \textbf{豆角爸爸回一句什么} \\
\midrule
量出来总多 1 厘米 & 把尺子边缘当成 0 & “手指点住 0，再放东西。” \\
\addlinespace[2pt]
从 3 量到 8 读成 8 厘米 & 不会用刻度差 & “数格子，或者 $8-3=5$。” \\
\addlinespace[2pt]
铅笔长 17 米 & 没有单位感 & “17 米比教室还长，铅笔有那么长？” \\
\addlinespace[2pt]
$1\ \text{米}=10\ \text{厘米}$ & 把“10 个 10 厘米”记成 10 & 让小豆角自己在尺上数到 100。 \\
\addlinespace[2pt]
量黑板只写“2 米” & 零头被丢掉了 & 必须写成“2 米 35 厘米”。 \\
\addlinespace[2pt]
画完线不标长度 & 书写习惯没建立 & 画完立刻写“X 厘米”。 \\
\addlinespace[2pt]
遇到坏尺子就卡住 & 不敢离开 0 & “起点 + 长度 = 终点”，随便挑个起点试。 \\
\bottomrule
\end{tabularx}
\end{center}

\subsection{抽查：随机开枪（10 题）}

\begin{enumerate}[label=\arabic*.,itemsep=4pt]
  \item 一支铅笔长 17 \blank[1.6cm]。\qquad 2. 一幢楼高 20 \blank[1.6cm]。
  \item 一枚硬币厚约 2 \blank[1.6cm]。
  \item $1\ \text{米}=$ \blank[1.6cm] 厘米。\qquad
  5. $300\ \text{厘米}=$ \blank[1.6cm] 米。
  \item $1\ \text{米}15\ \text{厘米}=$ \blank[1.6cm] 厘米。
  \item $148\ \text{厘米}=$ \blank[1cm] 米 \blank[1cm] 厘米。
  \item 一条线段从 2 画到 9，长 \blank[1.6cm] 厘米。
  \item 用坏尺子画 5 厘米，可以从 4 画到 \blank[1.6cm]。
  \item 小豆角的身高大约 \blank[1.6cm] 厘米。
\end{enumerate}

\begin{ansbox}
1. 厘米\quad 2. 米\quad 3. 厘米\quad 4. 100\quad 5. 3\quad 6. 115\quad
7. 1 米 48 厘米\quad 8. 7\quad 9. 9\quad 10. 按实际填（120--135）\\
\textbf{第 6、7 题错}：回任务 4，“米+厘米”再练 5 道。
\end{ansbox}

\subsection{课后野路子}

\begin{itemize}
  \item \textbf{门框上画身高线}：每月量一次，让小豆角自己读数、自己记日期。
  \item \textbf{周末量三样家具}：先估再量，记在附录 D 的表格里。
  \item \textbf{数步子算距离}：先量出小豆角一步多长，再数从客厅到卧室走几步，
        算一算大概几米（这就是乘法的伏笔）。
  \item \textbf{超市看包装}：米袋上写“5 kg”、水瓶上写“1.5 L”——
        告诉小豆角重量和容量也有单位，但这学期只学长度。
  \item \textbf{买鞋}：让小豆角自己量脚长得多少厘米，再去看鞋码表。
\end{itemize}
